\addtocontents{toc}{\protect\Needspace{3\baselineskip}}
\originalchapter{8.821 课程详讲：对称、参数、场与边界源}
本章对应第 17 讲的弦/SYM 部分和第 18 讲的通用字典. 先在 $\theta=0$、共形真空中讨论 $N$ 与 $g_{\rm YM}$ 两个参数. 本章把字典作为对偶猜想的具体陈述, 逐条解释它怎样检验或使用, 不以“两个表格相同”充当完整非微扰证明.

\originalsection{D3 上的场论为何是四维 \texorpdfstring{$\mathcal N=4$}{N=4}}
十维 Type IIB D3 有一个四维规范场、六个横向位置标量和相应费米伙伴. 对 $N$ 张膜, 全部为 $U(N)$ 的伴随矩阵. 四维一组无质量场的在壳数为规范场2个偏振加六个实标量, 共8个玻色自由度; 四个 Weyl 旋量共8个费米自由度, 总计16, 因此 $U(N)$ 总计 $16N^2$.

将每个矩阵分成无迹部分和单位矩阵部分, 因 $[I,M]=0$, 作用量中全部对易子相互作用不含单位矩阵部门. 得相互作用 $SU(N)$ 与自由质心 $U(1)$. 前者有 $16(N^2-1)$ 个领先在壳基本自由度, 后者16个. 重合膜整体的位置平移就在后者, 因而标准 AdS 喉部相互作用字典通常针对 $SU(N)$, 不能不说明地把那份自由质心也算成所有体模式的一部分.

$\mathcal N=4$ 表示四个 Weyl 形式的 Poincaré supercharge, 共16个实分量. 在共形真空, 还存在16个 superconformal charge, 共32. 六标量的旋转 $SO(6)$（或其自旋覆盖 $SU(4)$）是 R 对称, 连同 $SO(4,2)$ 共形群形成相应超共形结构. 原渐近平直 D3 保留一半超对称16个, 近视界 $\mathrm{AdS}_5\times S^5$ 则增强至32; 额外16个与场论的 superconformal charge 联系, 所以两次不同背景的计数不矛盾.

四维的 $g_{\rm YM}$ 经典无量纲, 本课程使用 $\mathcal N=4$ 理论的精确零 beta 函数性质, 在未给位置标量非零期望的原点真空中得到 CFT. “无量纲”本身并不能推出 beta 函数为零; 普通 Yang--Mills 是反例. 全阶超对称抵消需要另外的 Ward 恒等式/非重整化论证, 本章把它作为已知性质引入. 若进入分离膜的 Coulomb 真空, 标量期望值给出尺度, 共形对称会自发破缺, 不能继续把同一真空当作保持全部共形群.

\originalsection{标度维数怎样限制二点和三点函数}
初级算符在原点满足 $[K_\mu,\mathcal O(0)]=0$, 即由特殊共形生成元湮灭; 平移生成元对它的作用产生导数后裔. 对欧氏标量初级算符 $\mathcal O$, 采用坐标缩放 $x'=\lambda x$, $\mathcal O'(x')=\lambda^{-\Delta}\mathcal O(x)$ 的约定. $\Delta$ 是标度维数, 与工程维数在相互作用理论中可不同. 平移、旋转使二点函数只依赖 $r=|x_1-x_2|$, 尺度对称要求 $G_{12}\propto r^{-(\Delta_1+\Delta_2)}$.

反演 $x\to x/x^2$ 还让距离变成 $r'=r/(|x_1||x_2|)$. Primary 的权重必须分别与两点处的 Jacobian 缩放匹配; 若 $\Delta_1\ne\Delta_2$, 会多出依赖单个 $|x_i|$ 的因子, 与只依赖两点距离的形式不相容. 因此不同维数的标量初级算符二点为零. 在同维数的适当基中可写
\[
\langle\mathcal O_i(x_1)\mathcal O_j(x_2)\rangle
=\frac{C_{ij}}{|x_{12}|^{2\Delta}},\qquad
x_i\ne x_j.
\]
同维数退化算符需先对角化 $C_{ij}$, 不能不加基选择就把它写成单位阵. 接触项、张量算符结构和实时间的处方还需另说明.

三个标量初级算符的形式可设为 $|x_{12}|^{-a_{12}}|x_{23}|^{-a_{23}}|x_{31}|^{-a_{31}}$. 每一点的权给
\[
a_{12}+a_{31}=2\Delta_1,\quad
a_{12}+a_{23}=2\Delta_2,\quad
a_{23}+a_{31}=2\Delta_3.
\]
解得
\[
\langle\mathcal O_1\mathcal O_2\mathcal O_3\rangle
=\frac{C_{123}}{
|x_{12}|^{\Delta_1+\Delta_2-\Delta_3}
|x_{23}|^{\Delta_2+\Delta_3-\Delta_1}
|x_{31}|^{\Delta_3+\Delta_1-\Delta_2}}.
\]
对称性决定距离幂, 但不决定 $C_{ij},C_{123}$; 后面的体场计算正是为这些动力学系数提供工具. 四点以后出现交比, 共形对称留下的函数信息更多, 也不能直接称为全被对称性解出.

\begin{example}
三个维数都为2的标量, 三点距离因子为 $1/(|x_{12}|^2|x_{23}|^2|x_{31}|^2)$; 两点因子为 $|x_{12}|^{-4}$. 改变相互作用可以改变三点系数而仍保留完全同样的距离幂. 因而比对两个理论的共形群与尺度指数是必要检验, 还必须比对系数和更高点函数.
\end{example}

\originalsection{五球调和展开怎样产生五维场塔}
十维场可按 $S^5$ 调和函数展开. 例如一个简单最小耦合无质量标量,
\[
\Phi(x,z,\Omega)=\sum_{\ell,I}\phi_{\ell I}(x,z)Y_{\ell I}(\Omega),
\qquad-\nabla^2_{S^5}Y_{\ell I}
=\frac{\ell(\ell+4)}{R^2}Y_{\ell I}.
\]
球谐本征值可从六维欧氏调和齐次多项式 $r^\ell Y_\ell$ 读出: $\nabla_6^2$ 的径向部分贡献 $\ell(\ell+4)$, 要让总 Laplacian 为零, 角向部分就给相反本征值. $\ell=0$ 有一个常数模, $\ell=1$ 有六个模, 与六维坐标表示相同.

于是十维波动方程变为每个五维模式的 $(\nabla^2_{\mathrm{AdS}_5}-m_\ell^2)\phi_{\ell I}=0$, $m_\ell^2R^2=\ell(\ell+4)$. 更一般 supergravity 场还可混合, 并有张量/矢量球谐, 不能把这个简单标量质量公式误当作全部 IIB KK 谱. 量子数的匹配还要比对其 $SO(6)$ 表示和超对称多重态.

引力零模的有效耦合由作用量体积积分得到. 物理五球体积 $V_5=\pi^3R^5$, 因而
\[
\frac1{16\pi G_{10}}\int_{S^5}\dd^5\Omega\sqrt{g_{S^5}}
\int\dd^5x\sqrt{-g_5}R_5
=\frac1{16\pi G_5}\int\dd^5x\sqrt{-g_5}R_5,
\quad G_5=\frac{G_{10}}{\pi^3R^5}.
\]
五球等距 $SO(6)$ 的15个 Killing 向量还可用于度规的混合分量, 产生五维规范场. 它们与边界15个 R 对称流相联系; 全谱的详细降维是更大的工作, 此处只展示这一具体机制与简单标量塔.

\originalsection{参数字典与两个不同的近似极限}
沿本卷的规范场归一化,
\[
g_{\rm YM}^2=4\pi g_s,\qquad
\lambda=g_{\rm YM}^2N=\frac{R^4}{(\alpha')^2}.
\]
再用前章 $G_{10}$ 与球体积, 得
\[
\frac{G_{10}}{R^8}=\frac{\pi^4}{2N^2},\qquad
\frac{G_5}{R^3}=\frac{\pi}{2N^2},\qquad
\frac{\alpha'}{R^2}=\lambda^{-1/2}.
\]
左式与右式的量纲分别是 $L^8/L^8$、$L^3/L^3$ 和 $L^2/L^2$, 也说明把十维 $G$ 写成 $\alpha'^2$ 而非 $\alpha'^4$ 会立即不一致.

大 $N$ 压低相应体的量子引力环修正; 大 $\lambda$ 使体曲率半径远大于弦长, 压低有限弦长修正. 因而\textbf{边界强 't Hooft 耦合}可以有\textbf{弱曲率的经典引力描述}, 并不意味着边界弱耦合和体超引力同时适用. 为在微扰 IIB 中进一步保持 $g_s=\lambda/(4\pi N)\ll1$, 通常取 $1\ll\lambda\ll N$ 的大参数区间.

如果只取 $N\to\infty$、固定任意 $\lambda$, 弦环被压低, 但 $\alpha'/R^2$ 未必小. 这是 genus 首阶的弦理论极限, 仍可包含完整二维世界面量子涨落和有限弦长效应, 不等于把所有弦当成无振子的经典点粒子. $1/N^2$ 与 $1/\sqrt\lambda$ 是基本的修正控制参数; 某个可观测量的第一非零弦长修正可能由于超对称从更高幂开始, 不能把这张参数表当作每个量必有一项 $1/\sqrt\lambda$ 的精确级数.

\begin{example}
取 $N=1000,\lambda=100$, 得 $g_s\simeq0.00796$、$\alpha'/R^2=0.1$、$G_5/R^3\simeq1.57\times10^{-6}$. 这是两个不同修正量各自小的示例, 不是说有限数字100已消除一切修正. 若固定同样大 $\lambda$ 而让 $N$ 太小, $g_s$ 会变大, 弱弦环描述不再由这套参数保证.
\end{example}

\originalsection{径向尺度与边界能标的关系}
在 Poincaré AdS 的同一 $z$, 固有时与固有长度是 $\dd\tau=R\dd t/z$、$\dd\ell=R\dd x/z$. 故同一个局部物理过程映到边界单位后,
\[
E_{\rm boundary}=\frac RzE_{\rm local},\qquad
\ell_{\rm boundary}=\frac zR\ell_{\rm local}.
\]
对于固定局部尺度 $E_{\rm local}\sim1/R$, 边界能标是 $1/z$. 小 $z$ 对应高能和短距离, 大 $z$ 对应低能和长距离. 这个关系说明额外方向怎样记录场论尺度, 不是给边界再添加一条可以独立移动的普通空间轴.

截去 $z<\epsilon$ 区域, 对应边界紫外分辨率 $\Delta x\sim\epsilon$. 这个截止也去掉靠近 AdS 边界无限大的体固有距离和体积, 因而在\textbf{体几何体积}的意义下常叫红外截止. 这与“小 $z$ 对应\textbf{边界能标} 紫外”不矛盾, 两处红外/紫外指的量不同. 全局 AdS 的深处红移到有限值, 边界球面的尺度产生离散谱; 与平直边界的差别在 PS4.4 中完整求解.

\originalsection{全局对称为何对应体的规范场}
边界守恒流 $\partial_\mu J^\mu=0$ 以源 $a_\mu$ 耦合 $\int a_\mu J^\mu$. 变换 $a_\mu\to a_\mu+\partial_\mu\Lambda$ 时,
\[
\delta S=\int\partial_\mu\Lambda\,J^\mu
=-\int\Lambda\,\partial_\mu J^\mu+\text{边界项}=0
\]
在相应衰减或无边界条件下成立. 如果源是体场 $A_\mu(z,x)$ 的边界数据, 对这份源的冗余自然由体 $A_M\to A_M+\partial_M\widetilde\Lambda$ 扩展, $\widetilde\Lambda|_{\partial}=\Lambda$. 于是边界全局 $U(1)$ 流对应体局部 $U(1)$ 规范场.

体中在边界消失的规范参数是纯冗余; 在边界保留非平凡作用的适当渐近变换可产生边界全局荷. 类似地, 度规的保持边界条件的渐近等距变换对应边界共形对称. 不能把“快速消失的小规范变换”与“有非平凡边界荷的大规范变换”写成同一个概念.

边界度规源的变分是 $\delta S=\tfrac12\int\sqrt{g^{(0)}}T^{\mu\nu}\delta g^{(0)}_{\mu\nu}$. 体度规的 Fefferman--Graham 形式可写
\[
\dd s^2=\frac{R^2}{z^2}
[\dd z^2+g_{\mu\nu}(z,x)\dd x^\mu\dd x^\nu],
\qquad g_{\mu\nu}(z,x)\to g^{(0)}_{\mu\nu}(x).
\]
改变 $g^{(0)}$ 就是改变边界应力源, 所以 $T_{\mu\nu}$ 对应体度规涨落. 原始体 $G_{\mu\nu}$ 具有 $R^2/z^2$ 发散共形因子, 不能不去掉该因子就直接把 $G(0,x)$ 等同有限边界度规.

\originalsection{膨胀子源给出一个可核查例子}
令边界常值 $\Phi_\infty$ 决定 $g_s=\ee^{\Phi_\infty}$, 则 $g_{\rm YM}^{-2}=\ee^{-\Phi_\infty}/(4\pi)$. 若先把 SYM 欧氏作用量写成 $S=g_{\rm YM}^{-2}\int\mathcal L_0$, 在所选场归一化保持不变下,
\[
\delta S=-\frac1{g_{\rm YM}^2}
\int\dd^4x\,\delta\Phi_\infty(x)\mathcal L_0(x).
\]
所以膨胀子边界值是相应 Lagrangian 算符的源, 具体负号和系数由源与算符的定义吸收或保留, 不能无说明地把不同归一化的 $\mathcal L$ 全写成同一个算符. 与此同时, $SO(6)$ 流对应球面降维规范场, 应力对应体引力子.

通用字典因此是单迹共形初级算符与对应的单粒子体场联系, 后者的自旋、标度相关质量和其他量子数须匹配. 多迹算符通常对应多粒子态, 后裔对应同一表示的进一步态, 不需要对每一个形式局域积另添一套独立基本体场. 对一般有质量标量, 源也不是其未重标度的 $z=0$ 数值, 而是近边界某支幂律的系数; 下一章从波动方程实际推出该幂律和维数关系.

在有限耦合或有限 $N$ 的完整引力理论中,“$\mathrm{AdS}_5\times S^5$”指定的是渐近边界条件及通量部门, 并不要求每个激发态的整个内部度规都等于那一个真空乘积. 黑洞、波包和量子几何涨落同样属于相应部门. 这让字典能用于动力学状态, 而不是只比对两张固定背景的对称表.
